% English translation of checked_pages/esq_119_056.tex; source: cote 119, site 56, PDF page 57, typescript page 54.
% AI-translated by GPT-6 Astra (OpenAI, ChatGPT Pro mode).
\scanpage{119/56}
\mspage{54}

of Galois theory, probably in the very spirit of Galois.
The perpetual question ``and what is all this for ?'' seems to me to have
neither more nor less meaning in the case of the anabelian algebraic
geometry now coming into being than for Galois theory in
Galois's time (or even today, when the question is asked
by an overwhelmed student ...), and likewise for the comment
that generally accompanies it: ``all that is very general !''.

($^4$) It is therefore easy to conceive that a group such as $S\ell(2,\mathbb{Z})$, with
its ``arithmetic'' structure, is a veritable machine for constructing
``motivic'' representations of $\mathrm{Gal}(\bar{\mathbb{Q}}/\mathbb{Q})$ and of its
open subgroups, and that one thus obtains, at least in
principle, all motivic representations of weight 1,
or contained in a tensor product of such representations
(which already makes quite a bundle !). In 1981 I had begun
to experiment with this machine in a few particular cases,
obtaining various remarkable representations of $\mathrm{I}\!\Gamma$ in
groups $G(\mathbb{Z}^{\wedge})$, where $G$ is a group scheme (not necessarily
reductive) over $\mathbb{Z}$, starting from suitable homomorphisms
\[ S\ell(2,\mathbb{Z}) \longrightarrow G_o(\mathbb{Z}) , \]
where $G_o$ is a group scheme over $\mathbb{Z}$, and $G$ is constructed
from this as an extension of $G_o$ by a suitable group
scheme. In the ``tautological'' case $G_o = S\ell(2)_{\mathbb{Z}}$,
one finds for $G$ a remarkable extension of $G\ell(2)_{\mathbb{Z}}$ by a torus of
dimension 2, with a motivic representation that ``caps''\tnote{TN-056-01}{The quoted verb ``caps'' translates the figurative ``coiffe'', literally puts a cap on.} those
associated with the class fields of the extensions $\mathbb{Q}(i)$ and $\mathbb{Q}(j)$
(as chance would have it, the ``complex multiplication fields'' of the two
``anharmonic'' elliptic curves). There is a construction principle here
that seemed to me very general and very effective, but
I have not had (or taken) the leisure to work through its dévissage and follow it
to the end --- this is one of the many ``hot spots'' in the
program of foundations for anabelian algebraic geometry (or
``Galois theory'' in an enlarged version) that I propose to
develop. At present, and in an order of priority that is

\margin{56}
